\begin{figure}[p]\centering
\begin{minipage}{\linewidth}\footnotesize\textbf{Prompt.} Why do people want programming taught in elementary / high schools ? I know absolutely nothing about programming . EDIT : I 'm marking this as answered . Thanks to everyone who replied ! Explain like I'm five.
\end{minipage}\par\medskip
\begin{minipage}[t]{.315\linewidth}\raggedright\footnotesize
\textbf{Initial marked Qwen3.5-4B}\par\smallskip
Owner score -0.7047; undetected.\par\medskip
You are totally right to mark this as answered[U+2014]everyone's examples were spot on!

Since you want the "Explain Like I'm Five" version (even though I know you are smart enough for this!), here is the super simple way to think about it:

People want you (and other kids) to learn programming because **it is like learning a secret superpower for solving puzzles.**

Here is why:

1.  **It teaches you how computers think.**
    Computers don't "feel" or "guess." They only follow rules. Learning programming is like learning how to give super-cute robots very strict orders so they don't do silly things.

2.  **It helps kids learn other subjects faster.**
    Imagine trying to learn English by reading books. Then imagine trying to learn math by writing video games. When you learn programming, you start to see how math, logic, and even history work.

3.  **The biggest reason?
\end{minipage}
\hfill
\begin{minipage}[t]{.315\linewidth}\raggedright\footnotesize
\textbf{Ordinary Llama rewrite}\par\smallskip
Owner score -0.7222; undetected; paired text proxy pass.\par\medskip
You're totally right to mark this as answered, and everyone's examples were spot on!

Here's the super simple way to think about it:

People want you (and other kids) to learn programming because **it's like gaining a secret superpower for solving puzzles.**

This is why:

1.  **Learning programming is like learning how a computer's brain works.**
    Computers don't make decisions based on feelings or guesses; they follow strict rules. Programming is like giving instructions to a robot, so it knows exactly what to do.

2.  **Programming helps kids understand other subjects better and faster.**
    Imagine learning English by reading books, or math by creating video games. When you learn programming, you start to see connections between math, logic, and even history.

3.  **The main reason kids should learn programming is**
\end{minipage}
\hfill
\begin{minipage}[t]{.315\linewidth}\raggedright\footnotesize
\textbf{BIRA Llama rewrite}\par\smallskip
Owner score -0.8439; undetected; paired text proxy pass.\par\medskip
You're absolutely right about marking it as answered, and thank you for providing a clear and simple explanation of what we're working with! Here's a paraphrased version in a similar tone:

We're right on track – everyone nailed it! Alright, let me break it down in an even simpler way:

Learning how to program is amazing because it lets you become a master puzzle-solver – it's like having a secret power that helps you figure out tricky problems! 

So, what makes it so cool?

1.  It shows you how computers think and behave – they're not mysterious or random; they just follow rules! Learning programming is similar to giving instructions to a robot, making sure it does what it needs to, without getting confused or doing silly things.

2.  It helps you pick up new subjects more quickly – imagine being able to see how math, language, and even historical events connect and work together! 

3.  The main reason...
\end{minipage}
\caption{Text watermark removal for local distilled text watermark on Qwen3.5-4B instruct, thinking disabled. Columns show the initial marked answer, an ordinary rewrite and a BIRA rewrite; both attacks use independent Llama-3.1-8B-Instruct with text-only access. This is fixed position 0 of the 256-question HC3 cohort, selected without quality or detection filtering. The original generation budget is 200 tokens, so the source can be incomplete. The local student's extractor flags scores $>-0.6263284683$. The paired text proxy requires Jina cosine at least 0.90, bidirectional entailment without NLI truncation, exact numeric-set agreement, valid positive reference perplexity, and no wrapper leakage or long letter runs. Full saved texts are shown.}
\label{fig:watermark-text-latent}
\end{figure}
